% GENERATED FILE. Edit canonical lessons, then run npm run generate:books.
% canonical-source-hash: fnv1a64:7cf4291595f4ff11
% canonical-lessons: JA-C76-mayonaka, JA-C76-isshuukan, JA-C76-ichinichi, JA-C76-hannichi, JA-C76-tsuitachi

\chapter{Midnight, Whole week, One whole day, Half a day, First of the month}
\label{ch:ja-midnight-whole-week-one-whole-day-half-a-day-first-of-the-month}
\hlchaptermodality{76}

\begin{chapteropening}
\textbf{By the end of this chapter:} \emph{I can say midnight, a whole week, one whole day, half a day and the first of the month.}

Complete the last lesson of chapter 76: i can say midnight, a whole week, one whole day, half a day and the first of the month.
\end{chapteropening}

\section[\texorpdfstring{\ja{まよなか} (mayonaka)}{まよなか (mayonaka)}]{\ja{まよなか} (mayonaka) — midnight}

\label{lesson:JA-C76-mayonaka}



\begin{quote}
Before the new one: say the Japanese for the year before last, then the Japanese for oversleeping.
\end{quote}

\subsection*{You'll want to know: \ja{まよなか}}
\textbf{\ja{まよなか}} — \emph{mayonaka} — \textquotedblleft{}midnight\textquotedblright{}.

\begin{grammarlens}[title={What you've built}]
One more word for this chapter.
\end{grammarlens}

\subsection*{Your turn}
\begin{itemize}
\raggedright
  \item \emph{Say it:} \emph{mayonaka}
  \item \emph{Say it:} \emph{mayonaka}, once more
  \item \emph{Say it:} say \emph{mayonaka}
  \item \emph{Recall:} say the Japanese for the year before last, then the Japanese for oversleeping, then say \emph{mayonaka} again
\end{itemize}

\begin{tcolorbox}[breakable,colback=teal!4,colframe=teal!35!black,title={Before you move on}]
What does \ja{まよなか} mean? (\textquotedblleft{}Midnight\textquotedblright{}.) Say it once more.
\end{tcolorbox}
\section[\texorpdfstring{\ja{いっしゅうかん} (isshuukan)}{いっしゅうかん (isshuukan)}]{\ja{いっしゅうかん} (isshuukan) — a whole week}

\label{lesson:JA-C76-isshuukan}



\begin{quote}
Before the new one: say the Japanese for oversleeping, then the Japanese for midnight.
\end{quote}

\subsection*{You'll want to know: \ja{いっしゅうかん}}
\textbf{\ja{いっしゅうかん}} — \emph{isshuukan} — \textquotedblleft{}a whole week\textquotedblright{}.

\begin{grammarlens}[title={What you've built}]
One more word for this chapter.
\end{grammarlens}

\subsection*{Your turn}
\begin{itemize}
\raggedright
  \item \emph{Say it:} \emph{isshuukan}
  \item \emph{Say it:} \emph{isshuukan}, once more
  \item \emph{Say it:} say \emph{isshuukan}
  \item \emph{Recall:} say the Japanese for oversleeping, then the Japanese for midnight, then say \emph{isshuukan} again
\end{itemize}

\begin{tcolorbox}[breakable,colback=teal!4,colframe=teal!35!black,title={Before you move on}]
What does \ja{いっしゅうかん} mean? (\textquotedblleft{}A whole week\textquotedblright{}.) Say it once more.
\end{tcolorbox}
\section[\texorpdfstring{\ja{いちにち} (ichinichi)}{いちにち (ichinichi)}]{\ja{いちにち} (ichinichi) — one whole day}

\label{lesson:JA-C76-ichinichi}



\begin{quote}
Before the new one: say the Japanese for midnight, then the Japanese for a whole week.
\end{quote}

\subsection*{You'll want to know: \ja{いちにち}}
\textbf{\ja{いちにち}} — \emph{ichinichi} — \textquotedblleft{}one whole day\textquotedblright{}.

\begin{grammarlens}[title={What you've built}]
One more word for this chapter.
\end{grammarlens}

\subsection*{Your turn}
\begin{itemize}
\raggedright
  \item \emph{Say it:} \emph{ichinichi}
  \item \emph{Say it:} \emph{ichinichi}, once more
  \item \emph{Say it:} say \emph{ichinichi}
  \item \emph{Recall:} say the Japanese for midnight, then the Japanese for a whole week, then say \emph{ichinichi} again
\end{itemize}

\begin{tcolorbox}[breakable,colback=teal!4,colframe=teal!35!black,title={Before you move on}]
What does \ja{いちにち} mean? (\textquotedblleft{}One whole day\textquotedblright{}.) Say it once more.
\end{tcolorbox}
\section[\texorpdfstring{\ja{はんにち} (hannichi)}{はんにち (hannichi)}]{\ja{はんにち} (hannichi) — half a day}

\label{lesson:JA-C76-hannichi}



\begin{quote}
Before the new one: say the Japanese for a whole week, then the Japanese for one whole day.
\end{quote}

\subsection*{You'll want to know: \ja{はんにち}}
\textbf{\ja{はんにち}} — \emph{hannichi} — \textquotedblleft{}half a day\textquotedblright{}.

\begin{grammarlens}[title={What you've built}]
One more word for this chapter.
\end{grammarlens}

\subsection*{Your turn}
\begin{itemize}
\raggedright
  \item \emph{Say it:} \emph{hannichi}
  \item \emph{Say it:} \emph{hannichi}, once more
  \item \emph{Say it:} say \emph{hannichi}
  \item \emph{Recall:} say the Japanese for a whole week, then the Japanese for one whole day, then say \emph{hannichi} again
\end{itemize}

\begin{tcolorbox}[breakable,colback=teal!4,colframe=teal!35!black,title={Before you move on}]
What does \ja{はんにち} mean? (\textquotedblleft{}Half a day\textquotedblright{}.) Say it once more.
\end{tcolorbox}
\section[\texorpdfstring{\ja{ついたち} (tsuitachi)}{ついたち (tsuitachi)}]{\ja{ついたち} (tsuitachi) — the first of the month}

\label{lesson:JA-C76-tsuitachi}



\begin{quote}
Before the new one: say the Japanese for one whole day, then the Japanese for half a day.
\end{quote}

\subsection*{You'll want to know: \ja{ついたち}}
\textbf{\ja{ついたち}} — \emph{tsuitachi} — \textquotedblleft{}the first of the month\textquotedblright{}.

\begin{grammarlens}[title={What you've built}]
One more word for this chapter.
\end{grammarlens}

\subsection*{Your turn}
\begin{itemize}
\raggedright
  \item \emph{Say it:} \emph{tsuitachi}
  \item \emph{Say it:} \emph{tsuitachi}, once more
  \item \emph{Say it:} say \emph{tsuitachi}
  \item \emph{Recall:} say the Japanese for one whole day, then the Japanese for half a day, then say \emph{tsuitachi} again
\end{itemize}

\begin{tcolorbox}[breakable,colback=teal!4,colframe=teal!35!black,title={Before you move on}]
What does \ja{ついたち} mean? (\textquotedblleft{}The first of the month\textquotedblright{}.) Say it once more.
\end{tcolorbox}
